\section{System Participants, Roles, and Assumptions}

\subsection{Participants and Trust Boundaries}

We model four principal roles:
\begin{itemize}
    \item \textbf{Data owner} (e.g., doctor): uploads and manages private records.
    \item \textbf{Friend/authorized collaborator}: may access selected medium-level data.
    \item \textbf{General authenticated user}: can access only public outputs.
    \item \textbf{Hosting server}: trusted for enforcement of access control, typing, and audit policy.
\end{itemize}

Untrusted analysis programs are treated as attacker-controlled code submitted by otherwise valid users. Therefore, user authentication alone is insufficient; language-level flow control is required.

\subsection{Security Labels and Meaning}

Each stored object has a confidentiality label:
\begin{itemize}
    \item \texttt{private} $\equiv$ High: owner-only medical data,
    \item \texttt{friends} $\equiv$ Medium: owner and explicitly granted collaborators,
    \item \texttt{public} $\equiv$ Low: releasable artifacts and statistics.
\end{itemize}

Label semantics are enforced both at API level (who may call which command) and at program-analysis level (which values may influence which outputs).

\subsection{Attacker Model}

We assume a Dolev-Yao style network attacker outside trusted endpoints:
\begin{itemize}
    \item can intercept, replay, and reorder traffic,
    \item can initiate sessions and submit malicious programs,
    \item cannot break cryptographic primitives or compromise trusted server internals.
\end{itemize}

Additionally, we consider an \emph{insider misuse attacker}: a legitimate user attempting to extract private data through crafted code or indirect channels (for example branch behavior).

\subsection{Channel Abstraction and Justification}

The model uses abstract secure channels, corresponding to TLS-like endpoints between user and server. This abstraction is valid only under explicit guarantees:
\begin{itemize}
    \item \textbf{Confidentiality}: passive network observers cannot read plaintext payloads.
    \item \textbf{Integrity}: altered packets are detected and rejected.
    \item \textbf{Endpoint authentication}: the client talks to the intended server identity.
\end{itemize}

What we \emph{do not} assume from the channel abstraction:
\begin{itemize}
    \item no guarantee about application-level authorization logic,
    \item no guarantee about confidentiality once data is legitimately delivered to an endpoint,
    \item no prevention of leaks caused by accepted program semantics.
\end{itemize}

Hence channels are necessary but not sufficient. They protect data \emph{in transit}; the type system and API checks protect data \emph{in processing and storage}.

\subsection{Security Goals}

Our goals are stated as properties over traces and program runs:
\begin{itemize}
    \item \textbf{G1: Confidentiality/non-interference}. High inputs do not influence Low outputs, except via approved declassification.
    \item \textbf{G2: Controlled collaboration}. Medium data is visible only to owner and delegated collaborators.
    \item \textbf{G3: Integrity of sensitive state}. Low-trust inputs cannot overwrite High-integrity objects.
    \item \textbf{G4: Explicit accountability}. Every downgrade operation is auditable and policy-bound.
\end{itemize}

These goals are referenced throughout later chapters to ensure each design decision has a clear security purpose.

